\section{Notation and preliminary estimates}\label{sec:preliminaries}

All factor variables are positive integers unless another domain is specified.
We write $\Pplus(n)$ and $\Pminus(n)$ for the largest and least prime factors
of $n>1$, with $\Pplus(1)=1$ and $\Pminus(1)=\infty$.
An integer is $y$-smooth if $\Pplus(n)\le y$, and is rough above $y$ if
$\Pminus(n)>y$. We use $\varphi$ for Euler's function, $\tau$ for the divisor
function, and
\[
 e(t)=\exp(2\pi i t),\qquad
 (z)_j=z(z-1)\cdots(z-j+1),\qquad (z)_0=1.
\]
We write $\mu(n)=0$ if $n$ is divisible by a prime square, and
$\mu(n)=(-1)^r$ if $n$ is a product of $r$ distinct primes.
The symbol $\omega$ will count distinct prime divisors in the specified
prime groups; an unrestricted distinct-prime count will be identified when
it occurs. In particular, multiplicities in an integer do not increase
its group count.

Throughout the analytic argument,
\begin{equation}\label{eq:scales}
 L=\log x,\qquad T=\log L,\qquad W=\exp(\sqrt L),
\end{equation}
and $x$ tends to infinity. A dyad is an interval $[Y,2Y)$, or a fixed
constant enlargement of one; subinterval restrictions will be allowed
explicitly. We write $n\asymp Y$ when $n$ is bounded above and below by
fixed positive multiples of $Y$.

\begin{convention}[Parameters and uniformity]\label{conv:parameters}
Every constant is fixed before $x\to\infty$. A bound $L^{O(1)}$ has an
exponent independent of $x$. Its dependence on earlier fixed parameters
is permitted unless explicitly excluded. For a two-sided size statement
$Y=xL^{O(1)}$, we mean $|\log(Y/x)|=O(T)$.
Arbitrary logarithmic accuracy means a bound $O_A(L^{-A})$ for every
desired fixed $A$, with the auxiliary choices made in the specified order.

The geometric parameters used in the final sieve application are chosen
first. The required Type II accuracy is chosen next. Inside the analytic
argument the physical and ideal norm targets precede the comparison
kernels; the mark length $J$ is chosen after those kernels. The strength
of the character and Mellin discrepancy is chosen after the graph and
Fourier parameters, and the precision of the cofactor proxies is chosen
last. Each result below specifies the uniformity needed to respect this
order.
\end{convention}

All Hilbert spaces are complex. An operator defined by transitions acts
on a function at the target and sums or integrates its weighted values
at the input. Symmetrization of a list means the orthogonal projection
that averages all permutations of its coordinates.

\paragraph{A guide to recurring notation.}
The following table locates the objects used across sections. Their
full definitions and hypotheses appear at the indicated references.
\begin{center}
\small
\renewcommand{\arraystretch}{1.18}
\begin{tabular}{@{}p{.20\linewidth}p{.51\linewidth}p{.23\linewidth}@{}}
\toprule
Notation & Role & Definition \\
\midrule
$L,T,W$ & Logarithmic scales and the roughness cutoff. & \Cref{eq:scales} \\
$\mathcal P_g,V_g,\mu_g$
 & Prime groups, harmonic masses, and probability laws; the number
   of groups is $s\asymp T$. & \Cref{def:prime-groups} \\
$q,\omega_g,\omega$
 & Fixed damping parameter and distinct group-prime counts.
 & \Cref{def:prime-groups} \\
$J,\ell,M=J+\ell$
 & Fixed slot counts; $J$ is chosen after the comparison kernels.
 & \Cref{sec:transference} \\
$m$, later $m_{\rm probe}$
 & Fixed bound on probes per big group, independent of $J$.
 & \Cref{thm:ideal} \\
$W_{\mathbf t},W_\ell$
 & Marked divisor weights; in $W_\ell$ every group has $\ell$ marks.
 & \Cref{eq:marked-weight} \\
$n_*,F_0,G_0$
 & Group-free part and endpoint cores invariant under multiplication
   by group primes. & \Cref{sec:shifts}\par\Cref{eq:endpoint-form} \\
\bottomrule
\end{tabular}
\end{center}

\subsection{Classical prime estimates}

We use the prime number theorem with an error smaller than any fixed
negative power of the logarithm \cite[Corollary~39 and Exercise~40]{TaoSW},
and Mertens' estimates \cite[Theorems~15 and~26]{TaoMertens}
\begin{equation}\label{eq:mertens}
 \sum_{p\le y}\frac1p=\log\log y+O(1),\qquad
 \prod_{p\le y}\left(1-\frac1p\right)
 \sim\frac{e^{-\gamma_E}}{\log y}.
\end{equation}
Here and below sums or products indexed by $p$ are over primes. The
following forms of Siegel--Walfisz and the multiplicative large sieve
will be used. The former follows from its von Mangoldt formulation by
partial summation; see \cite[Exercise~64]{TaoSW}. For the latter see
\cite[Theorem~19.16]{MontgomeryVaughanII}.

\begin{theorem}[Siegel--Walfisz]\label{thm:sw}
For fixed $A,C>0$, uniformly for $(a,r)=1$ and
$1\le r\le(\log y)^C$,
\[
 \#\{p\le y:p\equiv a\pmod r\}
 =\frac{\Li(y)}{\varphi(r)}
  +O_{A,C}\!\left(\frac{y}{(\log y)^A}\right).
\]
The constants, which need not be effective, are independent of $y,a,r$.
\end{theorem}

\begin{theorem}[Multiplicative large sieve]\label{thm:large-sieve}
Let $I$ be a real interval of length $H$, let $a_n$ be complex numbers
supported on $I\cap\Z$, and let $R\ge1$. Then
\[
 \sum_{r\le R}\frac r{\varphi(r)}
 \sum_{\chi\bmod r}^{*}
 \left|\sum_{n\in I\cap\Z}a_n\chi(n)\right|^2
 \ll (H+1+R^2)\sum_n|a_n|^2,
\]
where the asterisk restricts the sum to primitive Dirichlet characters.
\end{theorem}

We will use absolute character estimates on intervals that may be very
short. The next consequence records precisely what is required; it
does not require a relative prime asymptotic on every such interval.

\begin{corollary}[Harmonic character estimates]\label{cor:harmonic-sw}
Fix $c,C,B,A>0$. If $\chi$ is nonprincipal modulo $r\le L^B$, then,
uniformly for real $|t|\le L^B$ and intervals
$I\subset[\exp(L^c),x^C]$,
\[
 \sum_{p\in I}\frac{\chi(p)p^{it}}p\ll_{c,C,B,A}L^{-A}.
\]
The same conclusion holds for a nonprincipal character induced from a
smaller modulus.
\end{corollary}

\begin{proof}
Writing $A_\chi(y)=\sum_{p\le y}\chi(p)$, sum
\Cref{thm:sw} against $\chi$ over reduced residue classes. The main
terms cancel, and the loss from the number of classes is at most $L^B$.
Because $\log y\ge L^c$, the permitted moduli are bounded by a fixed
power of $\log y$. Thus $A_\chi(y)\ll_K L^B y(\log y)^{-K}$ for every
fixed $K$. Partial summation on $I=(a,b]$ against $y^{-1+it}$ gives
boundary terms of size $O_K(L^B(\log a)^{-K})$ and an integral bounded by
\[
 C_KL^B(1+|t|)\int_a^b\frac{dy}{y(\log y)^K}
 \ll_K L^{2B}(\log a)^{1-K}.
\]
Choose $K$ sufficiently large. The argument is uniform in both
endpoints and also covers an empty interval. An induced nonprincipal
character is itself nonprincipal on the modulus on which it is used,
so the same argument applies.
\end{proof}

\subsection{Divisors, Dirichlet polynomials, and smooth separation}

\begin{lemma}[Fixed divisor moments]\label{lem:divisor-moments}
For each fixed nonnegative integer $k$,
\[
 \sum_{n\le Y}\frac{\tau(n)^k}{n}\ll_k(\log(2Y))^{C_k},
 \qquad
 \sum_{n\le Y}\tau(n)^k\ll_k Y(\log(2Y))^{C_k}
\]
for a constant $C_k$. Consequently, a convolution of a fixed number
of sequences bounded by fixed logarithmic powers has coefficient-square
sum $U L^{O(1)}$ on a dyad of size $U\le x^{O(1)}$.
If its coefficients include the reciprocal of the product index, this
bound is $L^{O(1)}/U$ instead.
\end{lemma}

\begin{proof}
Positivity and unique factorization bound the harmonic sum by
\[
 \prod_{p\le Y}\sum_{a\ge0}\frac{(a+1)^k}{p^a}
 =\prod_{p\le Y}\left(1+\frac{2^k}{p}+O_k(p^{-2})\right)
 \ll_k (\log(2Y))^{C_k},
\]
using \Cref{eq:mertens}. The counting bound follows by multiplying by
$Y$. If there are $b$ convolution factors, their coefficient at $n$
has absolute value at most a fixed logarithmic power times the number
of ordered $b$-factorizations of $n$. The latter is at most
$\tau(n)^{b-1}$: choose the first $b-1$ divisors, which determine the
last factor. Apply the counting moment bound with $k=2(b-1)$.
On a dyad the reciprocal index is comparable to $1/U$.
\end{proof}

The next coefficient-independent estimate is a weaker elementary form
of the Dirichlet-polynomial mean-value theorem; compare
\cite[Theorem~2 and Corollary~3]{MontgomeryVaughan1974}.

\begin{lemma}[Mean squares]\label{lem:dirichlet-mean}
If $P(t)=\sum_{n\le Y}c_n n^{it}$ and $I$ is an interval of length $H$,
then
\[
 \int_I|P(t)|^2\,dt
 \ll \bigl(H+Y\log(2Y)\bigr)\sum_n|c_n|^2.
\]
For a set $\mathcal T$ of real points at mutual distance at least one,
contained in an interval of length $H$, one also has
\begin{equation}\label{eq:discrete-mean}
 \sum_{t\in\mathcal T}|P(t)|^2
 \ll \bigl(H+1+Y\log(2Y)\bigr)(\log(2Y))^2
       \sum_n|c_n|^2.
\end{equation}
The constants are absolute, independently of any factorization used
to form the coefficients $c_n$.
\end{lemma}

\begin{proof}
Integration gives the diagonal $H\sum|c_n|^2$. For $n\ne m$, the
absolute value of the exponential integral is at most
$2/|\log(n/m)|\ll Y/|n-m|$. Apply
$2|c_nc_m|\le |c_n|^2+|c_m|^2$ and sum $1/|n-m|$ to obtain the first
bound. On the unit interval centered at $t$, the elementary
one-dimensional Sobolev inequality bounds $|P(t)|^2$ by a constant
times the integral of $|P|^2+|P'|^2$. These unit intervals have bounded
overlap. The derivative has coefficients $i(\log n)c_n$, so the first
bound applied twice gives \Cref{eq:discrete-mean}.
\end{proof}

We also use two elementary exponential-sum estimates in the following
forms; see \cite[Corollary~16.6 and Theorem~16.7]{MontgomeryVaughanII}.

\begin{lemma}[Derivative tests]\label{lem:derivative-tests}
Let $f$ be real-valued on an interval containing $H$ consecutive
integers.
\begin{enumerate}
\item If $f'$ is monotone and stays in
$[j+\lambda,j+1-\lambda]$ for some integer $j$ and
$0<\lambda\le1/2$, then $\sum_n e(f(n))\ll\lambda^{-1}$.
\item If $f$ is twice continuously differentiable and
$\lambda\le|f''|\le C\lambda$ throughout the interval, then
$\sum_n e(f(n))\ll_C H\lambda^{1/2}+\lambda^{-1/2}$.
\end{enumerate}
Both estimates hold on every subinterval on which the stated
hypotheses hold.
\end{lemma}

\begin{lemma}[Fourier separation]\label{lem:fourier-separation}
Let $F$ be smooth, supported in a fixed bounded box in $\R^b$, where
$b$ is fixed. Suppose that for every nonnegative integer $j$ its
derivatives of order $j$ are bounded by $C_jL^{C(j+1)}$.
Then Fourier inversion separates the variables with integrated
absolute coefficient mass $L^{O(1)}$.
There is a fixed $C'>0$, depending on $b,C$, such that restricting
each Fourier frequency to absolute value at most $L^{C'}$ leaves
an error $O_A(L^{-A})$ for every fixed $A$.
The exponent $C'$ may be fixed before the desired $A$.
\end{lemma}

\begin{proof}
Repeated integration by parts gives a bound for $\widehat F(\xi)$
by a derivative norm times an arbitrary fixed negative power of
$1+|\xi|$. Taking more than $b$ derivatives first makes this bound
integrable, with a fixed power of $L$. For the tail, take $j>b$
derivatives and integrate outside $|\xi|\ge L^{C'}$. The resulting
bound is at most a constant depending on $j$ times
\[
 L^{C(j+1)-C'(j-b)}.
\]
Fix $C'>C$ sufficiently large for the chosen Fourier convention and
box. Increasing $j$ then supplies any prescribed negative power of
$L$. Fourier inversion expresses the integrand as a product of
one-variable phases. Applying this in normalized logarithmic
coordinates gives the corresponding Mellin separation on fixed
dyads. Fixed-dimensional smooth cutoffs localized at logarithmic
precision are included in the derivative hypothesis.
\end{proof}

\subsection{An elementary weighted sieve}

We record a form with explicit coefficient and level bounds. This
will be used both for congruence counts and for oscillatory sums over
rough integers.
The construction is a version of the Brun--Hooley sieve: the product
of block upper bounds and its one-block correction are the inequalities
of Ford and Halberstam~\cite[Section~1, Lemma~1]{FordHalberstam2000}.
We prove the form needed here, including its coefficient and level bounds.

\begin{lemma}[Block sieve]\label{lem:sieve}
Consider a finite set of objects with nonnegative weights. For some
primes $p\le z$, let a bad condition be specified. If $d$ is a
squarefree product of these primes, suppose the weight of the objects
on which all conditions at $p\mid d$ hold is
\[
 Xg(d)+r(d),
\]
including $d=1$, where $X\ge0$, $g$ is multiplicative, and
\begin{equation}\label{eq:sieve-density-hypotheses}
 0\le g(p)\le1-\eta_0,\qquad
 \sum_{w<p\le w^2}g(p)\le C\quad(w>1)
\end{equation}
for fixed $\eta_0>0$ and $C$.
For every sufficiently large even integer $h$, the weight $S$ of
objects avoiding all bad conditions satisfies
\begin{equation}\label{eq:block-sieve}
 S=X\prod_{p\le z}(1-g(p))\bigl(1+O(e^{-h})\bigr)
   +O\!\left(\sum_{d\le z^{4h+2}}|r(d)|\right).
\end{equation}
Only the indicated primes and their squarefree products are used.
The implied constants depend on $\eta_0,C$, uniformly also when $h$
grows. For the fixed even choice $h=2$, there is an upper bound by
a constant times $X\prod(1-g(p))$ plus the same remainder sum.

More precisely, the proof supplies polynomials $U$ and $V$ in the
indicators of the bad conditions such that
\[
 V\le\ind_{\text{no bad condition}}\le U,\qquad U\ge0.
\]
Both have coefficients of absolute value at most one, supported on
squarefree $d\le z^{4h+2}$. The nonnegative overcount
$U-\ind_{\text{no bad condition}}$ is dominated by the pointwise
nonnegative polynomial $U-V$, whose absolute coefficients are also
at most one and which has the same level bound.
\end{lemma}

\begin{proof}
Partition the indicated primes into blocks
\[
 \mathcal B_j=\{p:z^{2^{-j-1}}<p\le z^{2^{-j}}\},\qquad j\ge0.
\]
Only finitely many blocks are nonempty. Set $h_j=(j+1)h$, which is
even. In block $j$, write $I_j$ for avoidance of all its bad conditions,
$U_j$ for the inclusion--exclusion polynomial through degree $h_j$,
and $E_j$ for the sum of all monomials of degree $h_j+1$.
If exactly $t$ bad conditions hold in the block, then
\[
 U_j=\sum_{a=0}^{h_j}(-1)^a\binom ta
 =\begin{cases}
  1,&t=0,\\
  \binom{t-1}{h_j},&t\ge1.
 \end{cases}
\]
Here a binomial coefficient is zero when its lower index exceeds
its nonnegative upper index. It follows that $U_j\ge I_j\ge0$ and
$U_j-I_j\le E_j$. Telescoping a product of nonnegative factors gives
\[
 0\le\prod_jU_j-\prod_jI_j
 \le\sum_jE_j\prod_{k\ne j}U_k.
\]
The upper and lower polynomials are therefore
\begin{equation}\label{eq:sieve-polynomials}
 U=\prod_jU_j,\qquad
 V=U-\sum_jE_j\prod_{k\ne j}U_k.
\end{equation}
In $U$, every block has degree at most $h_j$. A monomial from
correction $j$ has degree exactly $h_j+1$ in block $j$ and at most
$h_k$ elsewhere. These supports are disjoint across $j$ and disjoint
from the support of $U$. Thus the absolute coefficients in $U,V$
are at most one. The same disjointness gives this bound for $U-V$.

The logarithm of a product in $U$, divided by $\log z$, is at most
\[
 \sum_{j\ge0}h(j+1)2^{-j}=4h.
\]
A correction adds at most one more unit to this upper bound. Hence
the stated level $z^{4h+2}$ is valid for all the polynomials above.

Now give the bad conditions independent probabilities $g(p)$.
Their block avoidance probability $a_j=\E I_j$ obeys
\[
 a_j=\prod_{p\in\mathcal B_j}(1-g(p))
 \ge\exp(-C/\eta_0)=:a_*>0,
\]
because $-\log(1-u)\le u/\eta_0$ for $0\le u\le1-\eta_0$.
Also
\[
 e_j:=\E E_j\le\frac{C^{h_j+1}}{(h_j+1)!},\qquad
 a_j\le\E U_j\le a_j+e_j.
\]
For all large even $h$,
\[
 R_h:=\sum_j e_j/a_j\ll e^{-h}.
\]
Indeed the factorial tails eventually decrease geometrically in
their index, and $h_j=(j+1)h$; the same sum is bounded by an absolute
constant when $h=2$. Independence between blocks yields
\[
 \prod_ja_j\le\E U\le\prod_ja_j\,e^{R_h},
 \qquad
 \E\sum_jE_j\prod_{k\ne j}U_k
 \le\prod_ja_j\,R_he^{R_h}.
\]
Consequently the expectations of both bounding polynomials differ
from $\prod a_j$ by $O(e^{-h})\prod a_j$ for large $h$.
For $h=2$, the upper expectation is at most a constant times
$\prod a_j$.

Evaluate the two polynomials in the original counting problem.
The main terms of their monomials are exactly their independent
expectations multiplied by $X$. The absolute remainder for either
polynomial is at most $\sum_{d\le z^{4h+2}}|r(d)|$, by the coefficient
bound. Squeezing $S$ between them proves the result.
\end{proof}

\begin{corollary}[A polynomial for rough integers]\label{lem:rough-sieve-polynomial}
For the ordinary bad conditions $p\mid n$ at primes $p\le z$, let
$U(n)=\sum_{d\mid n}\lambda_d$ be the upper polynomial in
\Cref{eq:sieve-polynomials}, and put $D=z^{4h+2}$. Then
\[
 U(n)\ge\ind_{\Pminus(n)>z}\ge0,\quad
 |\lambda_d|\le1,\quad
 \lambda_d=0\quad\text{unless $d\le D$ is squarefree},
\]
and
\begin{equation}\label{eq:rough-sieve-overcount}
 \sum_{n\in I\cap\N}
 \bigl(U(n)-\ind_{\Pminus(n)>z}\bigr)
 \ll |I|e^{-h}+D
\end{equation}
for every finite interval $I\subset[1,\infty)$ and all sufficiently
large even $h$. The bound is uniform in $h,z,I$. Moreover,
\begin{equation}\label{eq:rough-sieve-harmonic}
 \sum_d\frac{|\lambda_d|}{d}
 \le\prod_{p\le z}(1+p^{-1})\ll\log(2z).
\end{equation}
\end{corollary}

\begin{proof}
The pointwise claims and harmonic sum follow from the coefficient
and support bounds. For the overcount, use its pointwise upper
bound $U-V=\sum_jE_j\prod_{k\ne j}U_k$, which is nonnegative
pointwise and has absolute coefficients at most one.
In the interval $I$, the count of multiples of $d$ is $|I|/d+O(1)$.
The independent expectation of $U-V$ was bounded in the proof by
$O(e^{-h})\prod_{p\le z}(1-1/p)$; the sum of absolute remainders is
$O(D)$. Dropping the avoidance product proves
\Cref{eq:rough-sieve-overcount}. Finally,
$\prod(1+1/p)\le\prod(1-1/p)^{-1}\ll\log(2z)$ by Mertens.
\end{proof}
